\item \points{3e}

We will now implement maximum-likelihood estimation for Bayesian network parameters, which assumes every variable is observed. Implement the function \texttt{mle\_estimate}, which takes in a \texttt{BayesianNetwork} and a list of observations and returns a new \texttt{BayesianNetwork} with the parameters estimated by MLE.

You will have to support Laplace smoothing using the parameter \texttt{lambda\_param}.

\expect{
    An implementation of the \texttt{mle\_estimate} function in \texttt{submission.py}.

    You should first implement the helper function \texttt{accumulate\_assignment} to accumulate the counts for each assignment. You should also use the provided helper functions \texttt{init\_zero\_conditional\_probability\_tables} and \texttt{normalize\_counts} where applicable; you don't need to implement them.
}